\subsection{Sequential Controls: Synthetic Mechanism Check}
\label{sec:appendix:sequential_simulation}

The real-data class-matched sampler fixes class quotas and spreads cells across
donors within each class. To distinguish these operations, we implement four
conditions with common random cell priorities and nested development subject
sets: (A) uniform pooled-cell sampling; (B) restriction to the eight-subject
reference class set, reserving one cell per class before uniform filling;
(C) exact reference class quotas with uniform within-class sampling; and
(D) the same quotas with capacity-constrained donor round robin in randomized
donor order. At fixed total, exact class proportions and exact counts are
equivalent constraints, not separate ablations. This A condition is not the
legacy equal-cell-per-subject sampler. Sequential contrasts are order-dependent
descriptions, not causal mediation estimates.

We evaluate this implementation on synthetic data, not a semi-synthetic clinical
cohort. Each of 20 seeds generates six class centers $\mu_k\sim
\mathcal{N}(0,s^2I_{12})$ and 48 subjects with 400 cells each. Subject label
probabilities are $\operatorname{softmax}(\eta_p)$, where
$\eta_p\sim\mathcal{N}(0,m^2I_6)$, and features are
$x_{pi}=\mu_{y_{pi}}+b_p+\epsilon_{pi}$, with
$b_p\sim\mathcal{N}(0,a^2I_{12})$ and
$\epsilon_{pi}\sim\mathcal{N}(0,I_{12})$ independently. We vary
$a,m\in\{0,1.5\}$ and $s\in\{0.5,1.5\}$. The first 32 subjects form the
development pool; the remaining 16 are untouched test subjects. An unweighted
LR probe is refit with training-only standardization for each condition and
$P\in\{8,32\}$ at total 1,600 training cells. All conditions reuse the same
realization and test subjects within a seed.

\begin{table}[htbp]
\centering\small
\begin{tabular}{rrrccc}
\toprule
Offset & Mixture & Separation & A: Unmatched & C: Quota & D: Quota + donor \\
\midrule
0.0 & 0.0 & 0.5 & \shortstack{+0.001\\{}[-0.007,+0.007]} & \shortstack{+0.001\\{}[-0.007,+0.008]} & \shortstack{+0.001\\{}[-0.006,+0.007]} \\
0.0 & 0.0 & 1.5 & \shortstack{-0.000\\{}[-0.003,+0.001]} & \shortstack{-0.000\\{}[-0.003,+0.002]} & \shortstack{-0.000\\{}[-0.003,+0.001]} \\
0.0 & 1.5 & 0.5 & \shortstack{+0.024\\{}[-0.005,+0.075]} & \shortstack{-0.001\\{}[-0.012,+0.010]} & \shortstack{-0.002\\{}[-0.012,+0.007]} \\
0.0 & 1.5 & 1.5 & \shortstack{+0.002\\{}[-0.005,+0.009]} & \shortstack{+0.000\\{}[-0.005,+0.003]} & \shortstack{-0.000\\{}[-0.003,+0.005]} \\
1.5 & 0.0 & 0.5 & \shortstack{+0.028\\{}[-0.028,+0.103]} & \shortstack{+0.028\\{}[-0.023,+0.103]} & \shortstack{+0.031\\{}[-0.015,+0.104]} \\
1.5 & 0.0 & 1.5 & \shortstack{+0.043\\{}[-0.005,+0.114]} & \shortstack{+0.042\\{}[-0.006,+0.113]} & \shortstack{+0.042\\{}[-0.009,+0.115]} \\
1.5 & 1.5 & 0.5 & \shortstack{+0.027\\{}[-0.051,+0.133]} & \shortstack{+0.016\\{}[-0.059,+0.115]} & \shortstack{+0.025\\{}[-0.039,+0.103]} \\
1.5 & 1.5 & 1.5 & \shortstack{+0.040\\{}[-0.029,+0.130]} & \shortstack{+0.029\\{}[-0.054,+0.114]} & \shortstack{+0.048\\{}[-0.008,+0.136]} \\
\bottomrule
\end{tabular}
\caption{Synthetic mechanism check: subject-balanced Macro-F1 difference between 32 and 8 training subjects at fixed total 1,600 cells. Entries are means across 20 independent seeds with empirical 2.5th--97.5th seed percentiles, not confidence intervals of the mean. Condition B equals A in every run because support is already complete and the support floor does not alter selection. All eight parameter settings are reported.}
\label{tab:sequential_simulation}
\end{table}


With exchangeable subjects ($a=m=0$), mean effects are near zero. With label
mixture variation alone and $s=0.5$, the unmatched mean effect is $+0.0245$,
compared with $-0.0011$ after quota matching. Conversely, with subject offsets
alone and $s=0.5$, quota matching retains a mean effect of $+0.0280$, compared
with unmatched $+0.0281$. Thus matching does not invariably eliminate a breadth
benefit in this generator. Seed variability is substantial; these results do
not establish interval coverage, a universal correction, or the mechanism in
APT, COMBAT, or OneK1K. Support is complete in these simulations, so they do not
test missing-class recovery. Separate real-cohort A--D results are not included
in this version; the existing class-matched tables must not be interpreted as
that decomposition.

\subsection{Known-Null and Independent-Target Validation}
\label{sec:appendix:known_target}

Known generator parameters are not themselves a known learning effect.
We therefore define the target explicitly as
\begin{equation}
\theta_C=\mathbb{E}\!\left[S_C(32,Q_8)-S_C(8,Q_8)\right],
\end{equation}
where $S$ is subject-balanced pooled six-subtype Macro-F1 on untouched test
subjects, and $Q_8$ is the development-only eight-subject class-quota reference.
The expectation integrates over generated cohorts and sampling. This
overlap-quota policy target excludes gains from acquiring labels absent at
$P=8$; it is not the total practical value of recruiting more subjects.

We generate a hierarchy of three lineages with two sibling subtypes each in
eight dimensions. Fixed class centers are $\mu_k=2e_{\lfloor k/2\rfloor}
+(-1)^{k+1}s e_{3+\lfloor k/2\rfloor}$ for $k=0,\ldots,5$ (zero-indexed axes).
Features follow $x_{pi}=\mu_{y_{pi}}+b_p+u_{p,g(i)}+v_{h(p)}+\epsilon_{pi}$,
with independent isotropic Gaussian effects: donor standard deviation $a$,
ten-cell block standard deviation $c$, four-donor batch standard deviation
$b$, and residual standard deviation 1. Within-subject correlation can thus
arise at donor or cell-block level; these covariance components should not be
interpreted as independent empirical estimates of biological heterogeneity.
Subject label probabilities are softmax-normal with logit standard deviation
1.5, and the sixth subtype is available in only 10\% of subjects. Each cohort
contains 32 development and 40 test subjects, with 120 cells per subject.
The training total is 480 at both $P=8$ and $P=32$; preprocessing and LR fitting
follow the preceding synthetic probe. Train and test subjects and batches are
disjoint. All tasks retain the fixed six-class scoring ontology.

The base setting is $(a,c,b,s)=(0,0,0,0.8)$. Conditional on the label-only
quota reference, training features in each class are iid with the same
distribution for both subject budgets. Consequently, the C fitted-model
distribution, and hence expected test score, is identical: $\theta_C=0$
exactly. This argument does not require balanced natural class proportions or
complete support. Additional settings change only $a$ to 0.5 or 1.5; stress
settings start from $a=0.5$ and set $c=0.7$, $s=0.3$, or $b=1$, respectively.
In the batch setting, increased subject breadth also increases batch coverage;
the target remains a policy contrast rather than a separate donor effect.

For every non-null setting we estimate $\theta_C$ using 512 independent
reference cohorts, with a direct uniform-without-replacement index sampler
separate from the evaluated priority-sampling implementation. Another 100
independent cohorts evaluate A/B/C using paired subject pools, cell priorities,
and test sets. Reference and evaluation random-seed ranges are disjoint.
We report reference Monte Carlo standard errors (MCSE), estimator bias, and
RMSE. Normal Monte Carlo intervals for bias combine evaluation-mean variance
and independent reference-mean variance; these are not patient-bootstrap
coverage measurements or multiplicity-adjusted tests.

\begin{table}[htbp]
\centering\small
\begin{tabular}{lrrrr}
\toprule
Setting & Target (MCSE) & A: mean / RMSE & B: mean / RMSE & C: mean / RMSE \\
\midrule
Known null & 0 (exact) & +0.0202 / 0.0354 & +0.0119 / 0.0249 & +0.0004 / 0.0109 \\
Mild donor shift & 0.0110 (0.0009) & +0.0278 / 0.0341 & +0.0221 / 0.0264 & +0.0102 / 0.0191 \\
Strong donor shift & 0.0135 (0.0015) & +0.0332 / 0.0412 & +0.0316 / 0.0392 & +0.0168 / 0.0358 \\
Correlated cells & 0.0108 (0.0008) & +0.0301 / 0.0334 & +0.0264 / 0.0290 & +0.0102 / 0.0199 \\
Weak subtypes & 0.0049 (0.0010) & +0.0247 / 0.0361 & +0.0235 / 0.0346 & +0.0016 / 0.0222 \\
Batch shift & 0.0181 (0.0017) & +0.0318 / 0.0463 & +0.0299 / 0.0438 & +0.0192 / 0.0429 \\
\bottomrule
\end{tabular}
\caption{Known-null and independent Monte Carlo validation of the overlap-quota target. A, B, and C denote unmatched, support-matched, and exact-quota sampling. Non-null targets use 512 independent reference cohorts; estimator means and RMSE use 100 separate evaluation cohorts. MCSE denotes reference Monte Carlo standard error. The target is the expected effect of the C allocation policy, not unrestricted subject recruitment. All six settings are reported; RMSE is measured relative to the exact null or estimated reference target.}
\label{tab:scaling_ground_truth}
\end{table}


Under the exact null, A and B yield mean effects $+0.0202$ and $+0.0119$,
whereas C yields $+0.0004$ (bias interval $[-0.0017,0.0026]$). The independent
reference check is $+0.00043$ with MCSE $0.00053$, consistent with the analytic
null. Mild and strong donor shifts have reference effects $+0.0110$ and
$+0.0135$; C estimates $+0.0102$ and $+0.0168$, respectively, rather than
mechanically forcing both to zero. Higher shift variance does not guarantee
a proportionally larger learning benefit. Across all six settings, C has lower
RMSE against this target than A or B. Nevertheless, A/B differences can be
legitimate benefits of their own allocation policies: their error relative to
$\theta_C$ demonstrates estimand mismatch, not that all observed breadth gains
are spurious. Non-null targets remain Monte Carlo approximations. This
validation neither identifies real-cohort causal mediation nor establishes
coverage of the joint uncertainty procedure used on real data.
